\section{Teorema cinese del resto}
Preso un sistema di \textbf{congruenze} del tipo
\begin{equation*}
  \begin{cases}
    x \equiv a_1 \mod n_1 \\
    x \equiv a_2 \mod n_2 \\
    x \equiv a_3 \mod n_3 \\
  \end{cases}
\end{equation*}
dove \(a_i\) sono numeri interi ed i \(n_i\) sono a due e due \textbf{coprimi} 
(ovvero non sono multipli tra di loro), il teorema assicura che
il sistema ammette soluzioni. \\
Vorrei precisare che il teorema ci da una condizione \textbf{sufficciente}, non
necessaria, questo vuol dire che un sistema dove i \(n_i\) \underline{non} sono
coprimi tra di loro potrebbe comunque ammettere soluzioni. 